\documentclass{article}
\usepackage{amsmath,amssymb,hyperref}
\title{A virtual empty cut from staggered camera coverage}
\author{}
\date{26 November 2026}
\begin{document}
\maketitle

\paragraph{Pins.}
A possible object moves in a convex one-dimensional domain with speed
at most $V$. Each qualified negative rules out its specified closed
region at the \emph{exposure} time. The reports have authenticated
coverage assumed in the model, registered coordinates, and separate
receipt times. They contain no object label. Single-point propagation
is exact in this plant; one must reconstruct coupled support afresh for
multiple discs with collision, appearance or other dependencies.
No camera hardware or independent selection is used.

\paragraph{Sharp two-view certificate.}
If region $R_i$ is empty at $t_i$, a position at a target cut $\tau$
is ruled out when its $V|t_i-\tau|$ reach ball lies wholly inside $R_i$.
This entails erosion, not interpreting two arbitrary negative images
as one exposure. Take $[-w,w]$ as the target; the left view
$[-w-m,h/2]$ appears at $\tau-\delta$ and the right view
$[-h/2,w+m]$ at $\tau+\delta$. With $m\ge V\delta$, their eroded
sets cover the target if and only if $h\ge2V\delta$, within this
symmetric two-view layout. Equality closes the gap if the coverage is
closed. For smaller overlap the admissible trajectory
$x(t)=-V(t-\tau)$ avoids each region at its own exposure and occupies
the band center at the target cut. A stationary object outside both can
produce the same negative reports. Enough domain/preparation space is
required; the lower bound does not apply to a physically blocked path.

The minimal synthetic raster depicts a passive reflectance spot behind
each aperture. When $V=1$, exposures are $\pm1/2$ and overlap is
$1/2$, the moving and stationary histories produce identical all-zero
payloads of the same size. At overlap 1 the moving trace crosses closed
visible boundaries and is detected. For unequal exposure distances
from $\tau$ the sum of the two erosion radii replaces $2V\delta$.
If a historical cut is synthesized at time 0 between observed anchors
$-1,1$, it can provide the missing negative for the earlier no-swap
criterion (maximum gap $1\le 6/5$ in the replay). Both component frames
must first have actually reached the deciding actor. The same virtual
cut delivered after 1 cannot authorize the earlier act.

\paragraph{A strictly stronger same-information comparator.}
For an ordered list of valid captures, propagate the feasible values
by interval dilation $V\Delta$, intersect the physical domain and
subtract the observed empty interval. A reverse pass from later
exposures uses the same operation. At $\tau$, intersect the two:
a shared point suffices to glue the prefix and suffix because the
dynamics have no internal state beyond position. Straight interpolation
is admissible between any compatible pair. If global feasibility itself
is empty, report inconsistent evidence rather than prove safety by
empty support. All open/closed endpoints matter.
In particular, let $V=1$, negatives $[-1,2]$ at 0 and $[-1,1/2]$
at 1, query $x=0$ at 2. The first alone allows a left predecessor
and the second alone a right predecessor. On a path to 0,
$x(0)<-1$ and $x(1)>1/2$ must both hold, which violates unit speed
on the first interval. Independent erosion misses this exclusion.

The separate exact reference chooses a left or right open component
at each negative, then checks all joint timing inequalities by
strict difference constraints. A negative cycle, or a zero total
cycle containing a strict edge, is exactly a contradiction. In a
finite graph with no such cycle, one can shrink each strict bound
by a sufficiently small positive amount without introducing any
negative cycle. A shortest-path potential then solves the system.
This explains why a fixed-grid witness search is not the reference.

\paragraph{Clock worlds determine which proofs compose.}
Consider reported times $-1,1$, two base regions $[-6,1]$ and
$[-1,6]$, uncertain offsets $\pm1/2$, $V=1$, target band $[-1,1]$, and time 0.
If the two captures share a common offset, their relative separation
stays 2. Their individually eroded edges can move with that common
offset yet cover the whole target in each admitted world. Replacing
their uncertainty by independent offsets incorrectly inserts a case
with separation 3. That case is a genuine feasible center crossing
if those independent clocks are in fact admitted. Eroding each view
for its own worst offset also fails to certify the known common-clock
case. Clock joins concern captures; receipt delays still apply at the
consumer. The replay shows the same common-clock proof unavailable
until receipt of the second report.

\paragraph{Evidence and limits.}
The exact code agrees with the independent constraint comparator
on 5625 cases and the simple erosion threshold on 45 symmetric
controls; four Lean implications cover endpoint and read-cut details.
\href{https://mediatum.ub.tum.de/doc/1575358/2n3lb03tprh28cbg1n19faj8w.Koschi_TIV2020_final.pdf}{Koschi--Althoff}
use more general set propagation for occlusion in traffic.
\href{https://georgejpappas.org/papers/J46.pdf}{Fainekos--Pappas}
study temporal robustness and sampling. The result here explains
the scope of one virtual interface and does not outperform exact
trajectory propagation or claim novelty of the underlying methods.
This is not the independent nonradio substitution, an imaging
performance estimate, or a claim about an already issued effect.
\end{document}
